\documentclass[a4paper]{report}
% Shared preamble for the English probability and statistics documents.
% Based on the original course preamble. Compile with pdfLaTeX.
\usepackage[top=1.5cm,bottom=2cm,left=2cm,right=2cm]{geometry}
\usepackage[T1]{fontenc}
\usepackage{lmodern}
\usepackage[utf8]{inputenc}
\usepackage[english]{babel}

\usepackage{amsmath}
\usepackage{amssymb}
\usepackage{amsthm}
\usepackage{mathrsfs}
\usepackage{mathtools}
\usepackage{etoolbox}
\usepackage{tikz}
\usetikzlibrary{shapes.multipart,arrows.meta}
\usepackage{xcolor}
\usepackage{enumerate}
\usepackage[shortlabels]{enumitem}
\usepackage{hyperref}
\usepackage{version}
\usepackage{bbold}
\usepackage{fontawesome5}
\usepackage{fancyhdr}
\lfoot{}
\usepackage{natbib}
\usepackage{graphicx}
\usepackage{pdfpages}
\usepackage{authblk}
    
\usepackage{blindtext}
\usepackage{multicol}
\usepackage{collcell}
\usepackage{wrapfig}
\usepackage{blkarray}
\usepackage{cancel}
\usepackage{appendix}
\usepackage{float}
\makeatletter
\input{tex/assets/slashbox.sty}
\makeatother
\hypersetup{unicode,colorlinks,citecolor=red,filecolor=black,
  linkcolor=red,urlcolor=cyan}
\setlength{\columnsep}{1cm}
\pagenumbering{gobble}

\def\numtext#1{%
  \ifcase#1\or one\or two\or three\or four\or five\or six\or seven\or eight\or nine\or ten\or
  eleven\or twelve\or thirteen\or fourteen\or fifteen\or sixteen\or seventeen\or eighteen\or nineteen\or twenty\or Lots\fi}

\newcommand{\Byte}[1]{%
  \def\nodecontents{}%
  \foreach[count=\l] \k in {#1} {\xappto{\nodecontents}{\noexpand\nodepart{\numtext{\l}}\k}}%
  \node[rectangle split,rectangle split parts=8,rectangle split horizontal,draw] at (2,2){\nodecontents};}

\newcommand{\D}{{\mathrm{d}}}
\newcommand{\cE}{{\mathcal{E}}}
\newcommand{\cG}{{\mathcal{G}}}
\newcommand{\cX}{{\mathcal{X}}}
\newcommand{\cY}{{\mathcal{Y}}}
\newcommand{\cB}{{\mathcal{B}}}
\newcommand{\cF}{{\mathcal{F}}}
\newcommand{\cO}{{\mathcal{O}}}
\newcommand{\cS}{{\mathcal{S}}}
\newcommand{\cT}{{\mathcal{T}}}
\newcommand{\given}{\, \big| \,}
\newcommand{\cL}{{\mathscr{L}}}
\newcommand{\cD}{{\mathcal{D}}}
\newcommand{\cA}{{\mathscr{A}}}
\newcommand{\cP}{{\mathscr{P}}}
\newcommand{\cR}{{\mathscr{R}}}
\newcommand{\A}{{\mathcal A}}
\newcommand{\one}{\mathbb{1}}
\newcommand{\mathbbm}[1]{\mathbb{#1}}

\providecommand{\C}{}
\renewcommand{\C}{\mathbb C}
\newcommand{\R}{\mathbb R}
\newcommand{\Q}{\mathbb Q}
\newcommand{\N}{\mathbb N}
\newcommand{\T}{\mathbb T}
\newcommand{\F}{\mathbb F}
\renewcommand{\P}{\mathbb P}
\newcommand{\Z}{\mathbb Z}
\newcommand{\E}{\mathbb {E}}
\newcommand{\Pbb}[1]{\P\left( #1\right)}
\newcommand{\Ebb}[1]{\E\left[ #1\right]}

\newcommand{\lb}{\left (}
\newcommand{\rb}{\right )}
\newcommand{\lbb}{\left [}
\newcommand{\rbb}{\right ]}
\newcommand{\labs}{\left |}
\newcommand{\rabs}{\right |}
\newcommand{\lbrb}[1]{\lb #1 \rb}
\newcommand{\lbbrbb}[1]{\lbb#1\rbb}
\newcommand{\labsrabs}[1]{\labs#1\rabs}
\newcommand{\lbbrb}[1]{\lbb#1\rb}
\newcommand{\lbrbb}[1]{\lb#1\rbb}
\newcommand{\langlerangle}[1]{\langle#1\rangle}
\newcommand{\lbcurly}{\{}
\newcommand{\rbcurly}{\}}
\newcommand{\curly}[1]{\lbcurly#1\rbcurly}
\newcommand{\abs}{\labsrabs}
\newcommand{\ind}[1]{1_{#1}}
\newcommand{\indep}{\perp \!\!\! \perp}
\renewcommand{\d}{\, \mathrm{d}}

\newtheorem{theorem}{Theorem}
\newtheorem{problemCnt}{Lemma}
\theoremstyle{definition}
\newtheorem{problem}[problemCnt]{Problem}
\newtheorem{problembook}[problemCnt]{\faBook \,  Problem}
\newtheorem{problempencil}[problemCnt]{\faPencil* \,  Problem}
\newtheorem{problemstar}[problemCnt]{\faStar\,  Problem}
\newtheorem{problemstarpencil}[problemCnt]{\faStar \, \faPencil* \,   Problem}
\newtheorem{probleminteresting}[problemCnt]{\faSmile[regular] \,   Problem}
\newtheorem{solution}{Solution}

\newcommand{\fwcell}[1]{\makebox[\arraystretch\normalbaselineskip]{$#1$}}
\newenvironment{enumeratei}{\begin{enumerate}[(i)]}{\end{enumerate}}
\newenvironment{enumeratea}{\begin{enumerate}[(a)]}{\end{enumerate}}
\newenvironment{enumeraten}{\begin{enumerate}[\upshape 1.]}{\end{enumerate}}

% Compatibility helpers needed by a few files in the combined archive.
\AtBeginDocument{\noindent{\footnotesize English translation prepared with ChatGPT (OpenAI, GPT-6), 26 October 2026, for translation of the Bulgarian text into English and LaTeX formatting.
The Bulgarian originals remain the reference; translation may contain errors.}\par\smallskip}
\newcommand{\MissingGraphic}[1]{%
  \begin{center}\fbox{\parbox{0.82\linewidth}{\centering Missing source image: \texttt{#1}}}\end{center}}



\setlength{\emergencystretch}{3em}
\allowdisplaybreaks
\begin{document}

\noindent Probability and Statistics, 01.12.2024 \hfill Test 1

\hrulefill

\begin{it}
The exact grading formula depends on the results. As an approximation, use 2 plus the points earned. Time allowed: 210 minutes. Good luck.

Some formulas: for $|x|<1$,
$\sum_{n=0}^\infty  x^n/n! = e^x$,
$\sum_{n=0}^\infty x^n = \frac{1}{1-x}$,
$\sum_{n=0}^\infty nx^n = \frac{x}{(1-x)^2}$,
$\sum_{n=0}^\infty n^2x^n = \frac{x(1+x)}{(1-x)^3}$,
$\sum_{n=0}^\infty x^n \binom{n+k}{n} = \frac{1}{(1-x) ^{k+1}}$.



\end{it}

\begin{problem} (1.75 points in total)
    \begin{enumerate}

    \item (0.25 points) Let \(\P(A)=0.4\), \(\P(B)=0.3\) and \(\P(\overline{A\cup B})=0.42\). Are \(A,B\) independent?
        \item (0.25 points) Five participants compete in a mathematics quiz. Each question has four possible answers, only one correct. Each participant answers five questions, choosing randomly. Find the probability that a given participant answers exactly three correctly.
        What is the probability that the correct answers were to questions 1, 3 and 5,
        given that exactly three answers were correct?


        \item (0.25 points) Let $X,Y,Z$
        be independent random variables. Prove that
        $Cov(X+Y,Z) = Cov(X,Z) + Cov(Y,Z)$.
        Does this remain true if $X,Y,Z$ are not
        independent?

\item (0.25 points) A fair coin is tossed ten times. Find the correlation between the number of heads in the first seven tosses and tails in the last seven tosses.

        \item (0.25 points) Ten people play a card-drawing game. Each draws one card from a 52-card deck without replacement. A player wins by drawing an ace. If anyone wins, the game ends; otherwise, a new round is drawn from the remaining cards. Find the probability that the game ends after exactly three rounds.

    \item (0.5 points)
Each face of a six-sided die is labelled with a uniformly
random number from $\{1,\dots,6\}$.
What is the probability
of rolling a six
with this die?
What is the expected number of rolls until the product of all results first becomes divisible by 6?
    \end{enumerate}

\end{problem}


\begin{problem} (1.25 points in total)
    Three friends play the following casino game:
  \begin{itemize}
        \item Each pays BGN 1 to enter;
        \item A fair coin is tossed until the first tail. The casino pays BGN 1 for each head;
        \item For each payment, the friends randomly choose who receives it.
  \end{itemize}
\begin{enumerate}
    \item (0.50 points) What is the probability
    that none of the three makes a loss?
    \item (0.75 points) What is each friend's expected profit? Also find its distribution and variance.
\end{enumerate}
\end{problem}

\begin{problem} (1.60 points in total)
Angel is a young Pokemon trainer who knows that ten species of flying Pokemon live near town.
He estimates that encountering a Pokemon takes about half an hour.
  \begin{enumerate}
    \item (0.50 points)
    What is the expected time needed to encounter all ten species? Assume the species are equally rare.
    \item (0.25 points) Let $X_i$ be the number of Pokemon seen until the first of species $i$. What is its distribution?
    Express the quantity in part 1 using the $X_i$.
    \item (0.10 points) Inclusion--exclusion applies
    in various contexts, including maxima and minima:
    $\max\{x_1,x_2\}=x_1+x_2-\min\{x_1,x_2\}$. Generalize this to $n$ numbers.
    \item (0.25 points) Let $Y_1,Y_2$
    be independent geometric random variables
    with parameters $p_1,p_2$. Find
    the distribution of $\min\{Y_1,Y_2\}$.
    \item (0.50 points)
    Angel has calculated the expectation in part 1, but realizes
    that assuming
    all ten species are equally rare
    is unrealistic. Can you answer part 1 if
    a random Pokemon's probability of being species $i$ is proportional to $i$? \textit{(For an explicit answer, solve the problem with four species.)}
  \end{enumerate}



\end{problem}


\begin{problem}\textit{(Bonus: 1.50 points)}
This problem introduces basic concepts of \emph{information theory}, developed by Shannon
in the mid-20th century, including \emph{entropy} and \emph{information}.

A code is a set of binary words,
for example $\{00,01,10,11\}$, intended
to transmit information through a channel, such as
left/right/up/down or any other four options.

All codes considered here must be \emph{prefix-free}.
This means no codeword is a prefix of another. Thus $\{0,00,001\}$ is not allowed, since $0$
is a prefix of $00$ and $001$, and $00$
is a prefix of $001$. Such a situation causes difficulties
for the receiver when decoding.


For example, the result of an experiment with 32 equally likely outcomes can be transmitted using all five-bit words.

Suppose four horses have winning probabilities $(1/2,1/4,1/8,1/8)$. We can encode the winner using
all two-bit words. But a code with a smaller expected number of bits is possible: encode horse 1 by $0$, horse 2 by $10$, horse 3 by $110$ and horse 4 by $111$.
\begin{itemize}
    \item (0.25 points) Is this code prefix-free? What is the expected number of bits needed to report the winner?

\end{itemize}



Entropy is a thermodynamic concept often described informally as a \emph{measure of disorder}, or, in probability, a \emph{measure of randomness}. If $X$ is discrete and $p_X(k):=\P(X=k)$, define its entropy by
\[
H(X) := \Ebb{-\log_2(p_X(X))} = - \sum_k \log_2(\P(X = k))\P(X=k).
\]
An engineering interpretation of entropy is \textit{the minimum expected number of bits needed to transmit the outcome of $X$}!

\begin{itemize}
    \item (0.25 points) Calculate the entropy in the last example and compare it with the expectation in the preceding part.
\end{itemize}

Let $X,Y$ be discrete random variables and
$p_{X,Y}(k,l):=\P(X=k,Y=l)$. Define
the joint entropy of $X,Y$
\[
H(X,Y) = - \Ebb{\log_2 ( P_{X,Y}(X,Y)} =  -\sum_{k,l} \log_2 \lbrb{P(X= k , Y=l)}P(X= k , Y=l)
\]
and the conditional entropy $H(X|Y)$ by
\[
H(X|Y) = \sum_{k,l} \P(X=k,Y= l) \log_2 \lbrb{\frac{\P(X=k,Y=l)}{\P(X=k)}}.
\]
$Y$ may provide information about $X$, thereby reducing its randomness. This change is called the \emph{mutual information of $X,Y$}:
\[
I(X,Y) = H(X) - H(X|Y).
\]

\begin{itemize}
    \item (0.25 points) Prove that
    $H(X,Y)=H(X)+H(Y|X)$ and that
    $I(X,Y) = I(Y,X)$.

    \item (0.25 points) Toss a fair coin, and let \(X\) count the tosses until the first head. Find \(H(X)\).
    \item (0.25 points) For this \(X\), propose an efficient sequence of yes/no questions of the form: Is \(X\) in \(S\)? Compare \(H(X)\) with the expected number of questions needed to determine \(X\).
\item (0.25 points) Let \(Y\) be discrete with values in $(0,1)$. Compare $H(Y)$ and $H(Z)$ when \(Z=2Y\), and when $Z=\cos(Y)$. Explain.

\end{itemize}


\end{problem}



\end{document}
